% generated by paper/make_tables.py -- do not edit
\begin{table*}[t]
\caption{Track B. Recovering the clean-signal NEWS2 band from sensor-corrupted
telemetry. All rows see the same corrupted observations. Missed-HIGH is the
fraction of HIGH observations reported as LOW.}
\label{tab:trackb}
\centering
\begin{tabular}{lrrrrr}
\toprule
Model & Accuracy & Macro-F1 & HIGH recall & Missed-HIGH & Mean cost \\
\midrule
news2\_rule\_on\_observed & 0.9005 & 0.8006 & 0.8847 & 0.0127 & 0.3165 \\
threshold\_alarms & 0.7352 & 0.6280 & 0.8520 & 0.0070 & 0.4461 \\
logistic\_regression & 0.8677 & 0.7866 & 0.8398 & 0.0023 & 0.2893 \\
random\_forest & 0.9005 & 0.8267 & 0.9013 & 0.0014 & 0.2074 \\
gradient\_boosting & 0.9260 & 0.8324 & 0.9199 & 0.0082 & 0.2651 \\
\bottomrule
\end{tabular}

\vspace{2pt}
\footnotesize Cohort: 120,000 readings from
30,000 simulated patients, split by patient.
Selection metric: \texttt{macro\_f1}. Macro-F1 selects \texttt{gradient\_boosting} while the asymmetric-cost criterion selects \texttt{random\_forest}; the F1-optimal model carries the higher rate of HIGH observations reported as LOW.
Cost weights are a stated assumption, not derived from outcome data.
\end{table*}
